% ============================================================
% CRC Lab — viva preparation sheet (read before the presentation).
% ============================================================
\documentclass[11pt,a4paper]{article}

\usepackage[margin=1in]{geometry}
\usepackage[T1]{fontenc}
\usepackage{lmodern}
\usepackage{xcolor}
\usepackage{enumitem}
\usepackage{titlesec}
\titleformat{\section}{\large\bfseries}{\thesection}{1em}{}
\titleformat{\subsection}{\normalsize\bfseries}{\thesubsection}{1em}{}

\begin{document}

\begin{center}
{\Large\bfseries CRC Lab --- Viva Preparation Sheet} \\[0.4em]
{\normalsize Read this once, then walk in. Total demo: about 2 minutes.}
\end{center}

\vspace{0.5em}
\hrule
\vspace{1em}

\section*{1. The one paragraph (memorize this)}

CRC appends $n-1$ zeros to the data, divides by the generator using XOR,
and the leftover is the CRC. We transmit data + CRC. The receiver divides
what arrives by the same generator: remainder zero means no error,
non-zero means a bit flipped. XOR replaces subtraction because binary
arithmetic has no borrows. If they ask nothing else, say this and offer to
type their example live.

\section*{2. Minute-by-minute demo}

\begin{enumerate}[leftmargin=1.5em]
  \item \textbf{Inputs (30s).} Default \texttt{HELLO} + \texttt{1011}.
  Read the numbers live: data bits, 3 appended zeros, remainder, frame.
  \item \textbf{Division (30s).} Press Play. Narrate: quotient on top,
  one staggered subtraction row per step, $\oplus$ where the leading 1
  was cancelled.
  \item \textbf{Error (30s).} Flip exactly one bit. Red bit, remainder
  non-zero, verdict ERROR. Say: one flipped bit, caught.
  \item \textbf{Reset + live input (30s).} Reset to VALID. Ask the
  examiner for a word, type it. Everything recomputes.
  \item \textbf{Report (15s).} Hit Export, show the PDF matches the
  screen: same inputs, every division step, same verdict.
\end{enumerate}

Golden rule: Reset before every flip demo. One flip, show, reset. Never
stack many flips on stage (Section~4 explains why).

\section*{3. Likely questions, one-line answers}

\begin{itemize}[leftmargin=1.5em]
  \item \textbf{Why append $n-1$ zeros?} To make room for the remainder;
  a remainder modulo an $n$-bit generator is always $n-1$ bits.
  \item \textbf{Why must the generator start and end with 1?} Leading 1
  sets its degree; trailing 1 means it catches all single-bit errors.
  \item \textbf{Why XOR, not subtraction?} In mod-2 arithmetic they are
  identical; XOR just has no borrows.
  \item \textbf{Does CRC correct errors?} No, detection only. It says
  \textit{that} something changed, never what the original was.
  \item \textbf{Text vs Bits mode?} Text converts each character to
  8-bit ASCII first; the division always runs on bits.
  \item \textbf{What did you build vs reuse?} Everything: the division
  engine, both UIs, the LaTeX exporter. Only \texttt{pdflatex} itself
  is external.
\end{itemize}

\section*{4. The trap: VALID with flipped bits}

If you flip several bits, their remainders can XOR to zero and cancel,
so the verdict reads VALID. This is genuine theory, verified during
development: single-bit errors are always caught; some multi-bit
patterns are exact multiples of the generator and slip through. If the
examiner spots it, say exactly that. Then reset and demo a single flip.

\section*{5. Day checklist}

\begin{itemize}[leftmargin=1.5em]
  \item \texttt{python3 report/server.py} running before they arrive;
  browser open at \texttt{http://127.0.0.1:8321/}.
  \item TUI binary tested once in their terminal size
  (\texttt{./target/release/crc-tui}): \texttt{d} edits message,
  \texttt{g} edits generator, \texttt{space} plays, \texttt{e} flips,
  \texttt{x} resets, \texttt{t} changes theme, \texttt{?} help.
  \item One sample PDF in hand or on disk.
  \item If stuck on any question: return to the memorized paragraph
  and offer a live example.
\end{itemize}

\end{document}
